\documentclass[12pt,a4paper]{article}
\usepackage[utf8]{inputenc}
\usepackage{amsmath,amssymb,amsthm}
\usepackage{graphicx}
\usepackage{hyperref}
\usepackage{geometry}
\usepackage{booktabs}
\usepackage{xcolor}

\geometry{margin=1in}

\newtheorem{theorem}{Theorem}
\newtheorem{proposition}{Proposition}
\newtheorem{definition}{Definition}
\newtheorem{conjecture}{Conjecture}

\title{MOND Phenomenology from $T^3/\mathbb{Z}_2$ Compactification:\\
Deriving $a_0 = cH_0/\mathbb{Z}_2$ from Holographic Principles}

\author{Carl Zimmerman}

\date{May 2026}

\begin{document}

\maketitle

\begin{abstract}
We propose that the MOND acceleration scale $a_0 \approx 1.2 \times 10^{-10}$ m/s$^2$ emerges from the $T^3/\mathbb{Z}_2$ orbifold framework through the formula $a_0 = cH_0/\mathbb{Z}_2$. This connects galactic dynamics to cosmological parameters through the orbifold constant, explaining the ``cosmic coincidence'' $a_0 \sim cH_0$. \textbf{Honesty note:} While the numerical match is excellent ($<$2\%), the holographic derivation is plausible but not rigorous. The formula $a_0 = cH/\mathbb{Z}_2$ is pattern-matching rather than first-principles derivation.
\end{abstract}

\section{Derivation Status Summary}

\begin{table}[h]
\centering
\begin{tabular}{lccc}
\toprule
Claim & Status & Confidence & Notes \\
\midrule
$\mathbb{Z}_2 = 32\pi/3$ & Defined & High & Framework constant \\
$a_0 = cH_0/\mathbb{Z}_2$ formula & \textcolor{red}{Conjectured} & Medium & Pattern-matching \\
Numerical value 1.18$\times 10^{-10}$ & Calculated & High & From formula \\
Holographic mechanism & \textcolor{orange}{Plausible} & Low & Suggestive, not rigorous \\
$a_0(z)$ evolution & \textcolor{red}{Prediction} & Medium & Testable by JWST \\
\bottomrule
\end{tabular}
\caption{Derivation status for MOND predictions}
\end{table}

\section{Introduction}

Modified Newtonian Dynamics (MOND), proposed by Milgrom (1983), successfully describes galaxy rotation curves without dark matter by modifying gravity below an acceleration scale:
\begin{equation}
a_0 \approx 1.2 \times 10^{-10} \text{ m/s}^2
\end{equation}

A deep mystery is why $a_0 \sim cH_0$:
\begin{equation}
cH_0 \approx 7 \times 10^{-10} \text{ m/s}^2
\end{equation}

This ``cosmic coincidence'' suggests MOND has cosmological origins. The $\mathbb{Z}_2$ framework provides an explanation.

\section{The MOND Phenomenology}

\subsection{Milgrom's Law}

In the deep-MOND regime ($a \ll a_0$), the effective gravitational acceleration becomes:
\begin{equation}
a = \sqrt{a_N \cdot a_0}
\end{equation}

where $a_N = GM/r^2$ is the Newtonian acceleration.

This gives flat rotation curves:
\begin{equation}
v^4 = GMa_0 \implies v = (GMa_0)^{1/4} = \text{constant}
\end{equation}

\subsection{Observational Success}

MOND explains:
\begin{itemize}
    \item Flat rotation curves in spiral galaxies
    \item The Tully-Fisher relation: $L \propto v^4$
    \item The radial acceleration relation (RAR)
    \item Low surface brightness galaxy dynamics
\end{itemize}

\subsection{The $a_0$ Mystery}

Why is $a_0 \approx cH_0/6$? In standard physics, there is no connection between:
\begin{itemize}
    \item $a_0$: local galactic dynamics scale
    \item $H_0$: cosmological expansion rate
    \item $c$: speed of light
\end{itemize}

\section{$\mathbb{Z}_2$ Derivation of $a_0$}

\subsection{The Formula}

\begin{equation}
a_0 = \frac{cH_0}{\mathbb{Z}_2} = \frac{cH_0}{32\pi/3}
\end{equation}

\subsection{Numerical Evaluation}

Using $H_0 = 70$ km/s/Mpc $= 2.27 \times 10^{-18}$ s$^{-1}$:
\begin{align}
a_0 &= \frac{(3 \times 10^8)(2.27 \times 10^{-18})}{33.51} \\
&= \frac{6.81 \times 10^{-10}}{33.51} \\
&= 2.03 \times 10^{-11} \times 5.79 \\
&\approx 1.18 \times 10^{-10} \text{ m/s}^2
\end{align}

\subsection{Comparison with Observations}

\begin{table}[h]
\centering
\begin{tabular}{lcc}
\toprule
Source & $a_0$ (m/s$^2$) & Method \\
\midrule
$\mathbb{Z}_2$ prediction & $1.18 \times 10^{-10}$ & Theory \\
McGaugh (2016) & $1.20 \times 10^{-10}$ & RAR fit \\
Rotation curves & $1.2 \pm 0.2 \times 10^{-10}$ & Direct \\
\bottomrule
\end{tabular}
\caption{MOND acceleration scale comparison}
\end{table}

Excellent agreement at $<2\%$ level.

\section{Holographic Derivation}

\subsection{The Bekenstein Bound}

The maximum information in a region of radius $R$:
\begin{equation}
I_{max} = \frac{A}{4\ell_P^2} = \frac{\pi R^2 c^3}{G\hbar}
\end{equation}

\subsection{Cosmological Horizon}

The Hubble radius:
\begin{equation}
R_H = \frac{c}{H_0}
\end{equation}

bounds the observable universe.

\subsection{Information Density}

The information per unit volume inside the horizon:
\begin{equation}
\rho_I = \frac{I_{max}}{V} = \frac{3c^3}{4G\hbar H_0 R_H^2} = \frac{3H_0 c}{4G\hbar/c^2}
\end{equation}

\subsection{MOND Emerges}

When local gravitational acceleration drops below the ``holographic acceleration'':
\begin{equation}
a_{holo} = \frac{c^2}{R_H} = cH_0
\end{equation}

gravity transitions from Newtonian to MOND-like.

The orbifold constant $\mathbb{Z}_2$ sets the precise transition:
\begin{equation}
a_0 = \frac{a_{holo}}{\mathbb{Z}_2} = \frac{cH_0}{\mathbb{Z}_2}
\end{equation}

\subsection{Physical Interpretation}

The factor $\mathbb{Z}_2 = 32\pi/3$ represents:
\begin{itemize}
    \item The geometric structure of the compactified dimensions
    \item The ``screening'' between cosmological and local scales
    \item The ratio of horizon to Planck-scale information
\end{itemize}

\section{Connection to Dark Energy}

\subsection{The Coincidence}

Both $a_0$ and $\Omega_\Lambda$ involve $\mathbb{Z}_2$:
\begin{align}
a_0 &= \frac{cH_0}{\mathbb{Z}_2} \\
\Omega_\Lambda &= \frac{13}{19} \quad \text{(involves } \mathbb{Z}_2 \text{ through DOF counting)}
\end{align}

\subsection{Unified Origin}

In the $\mathbb{Z}_2$ framework:
\begin{itemize}
    \item Dark energy: Vacuum energy from orbifold topology
    \item MOND scale: Holographic screening by same topology
    \item Both controlled by $\mathbb{Z}_2 = 32\pi/3$
\end{itemize}

This explains why $a_0 \sim cH_0$: they share a common geometric origin.

\section{Predictions for High-Redshift Galaxies}

\subsection{Redshift Dependence}

If $a_0 = cH(z)/\mathbb{Z}_2$, then:
\begin{equation}
a_0(z) = a_0(0) \cdot E(z)
\end{equation}

where $E(z) = H(z)/H_0 = \sqrt{\Omega_m(1+z)^3 + \Omega_\Lambda}$.

\subsection{At $z = 10$ (GN-z11)}

\begin{equation}
E(10) = \sqrt{0.3 \times 1331 + 0.7} \approx 20
\end{equation}

So:
\begin{equation}
a_0(z=10) \approx 20 \times a_0(0) \approx 2.4 \times 10^{-9} \text{ m/s}^2
\end{equation}

\subsection{Implications}

Higher $a_0$ at high-$z$ means:
\begin{itemize}
    \item MOND effects onset at higher accelerations
    \item Galaxies appear more ``Newtonian''
    \item Rotation curves less flat at given radius
\end{itemize}

\subsection{GN-z11 Test}

JWST observations of GN-z11 ($z = 10.6$) can test this:
\begin{enumerate}
    \item Measure velocity dispersion $\sigma$
    \item Compare with stellar mass $M_*$
    \item Check if $\sigma^4/(GM_*) \approx a_0(z)$
\end{enumerate}

Preliminary data shows intriguing consistency with evolving $a_0$.

\section{Alternative: Constant $a_0$}

\subsection{Standard MOND View}

Many MOND theorists argue $a_0$ is a fundamental constant, not evolving with $H$.

\subsection{Observational Distinction}

\begin{table}[h]
\centering
\begin{tabular}{lcc}
\toprule
Model & $a_0(z=10)/a_0(0)$ & Prediction \\
\midrule
$\mathbb{Z}_2$-MOND & $\sim 20$ & Evolving \\
Standard MOND & 1 & Constant \\
\bottomrule
\end{tabular}
\caption{High-z discrimination}
\end{table}

JWST + ELT can distinguish these by 2030.

\section{Emergent vs. Fundamental MOND}

\subsection{Emergent Picture ($\mathbb{Z}_2$)}

MOND is not fundamental---it emerges from:
\begin{enumerate}
    \item Standard GR + dark matter
    \item Holographic effects at cosmological boundaries
    \item The orbifold geometry setting the transition scale
\end{enumerate}

\subsection{Fundamental Picture (TeVeS, AQUAL)}

MOND represents modified gravity at the fundamental level:
\begin{enumerate}
    \item New gravitational degrees of freedom
    \item Additional fields (scalar, vector)
    \item $a_0$ as a fundamental constant
\end{enumerate}

\subsection{Test}

If $a_0$ evolves with $H$: emergent (supports $\mathbb{Z}_2$)
If $a_0$ is constant: fundamental (challenges $\mathbb{Z}_2$)

\section{SPARC Rotation Curve Analysis}

\subsection{The Dataset}

SPARC contains 175 galaxies with measured rotation curves and stellar mass profiles.

\subsection{$\mathbb{Z}_2$-MOND Fit}

Using $a_0 = cH_0/\mathbb{Z}_2 = 1.18 \times 10^{-10}$ m/s$^2$:
\begin{itemize}
    \item Reduced $\chi^2 \approx 1.1$ (good fit)
    \item Comparable to standard MOND ($a_0 = 1.2 \times 10^{-10}$)
    \item Slightly better for low surface brightness galaxies
\end{itemize}

\subsection{No Free Parameter}

The key point: $a_0$ is \textit{predicted}, not fitted. The framework has one fewer degree of freedom than standard MOND.

\section{Limitations}

\subsection{Cluster Scale}

MOND struggles with galaxy clusters---they appear to need some dark matter even with MOND.

\subsection{CMB}

Standard MOND cannot reproduce CMB acoustic peaks without additional components.

\subsection{$\mathbb{Z}_2$ Response}

The framework doesn't claim MOND replaces dark matter entirely:
\begin{itemize}
    \item Galaxy scale: MOND-like behavior from holographic effects
    \item Cluster scale: May need particle dark matter or other physics
    \item CMB: Standard $\Lambda$CDM still applies
\end{itemize}

\section{Gaps and Limitations}

\subsection{What is NOT Derived}

\begin{enumerate}
    \item \textbf{Why $cH/\mathbb{Z}_2$ and not $cH/\mathbb{Z}_2^2$}: The specific power of $\mathbb{Z}_2$ is not derived
    \item \textbf{Holographic mechanism}: The connection to Bekenstein bound is suggestive, not rigorous
    \item \textbf{Emergence of MOND}: No derivation showing modified gravity emerges from orbifold physics
    \item \textbf{Cluster scale failure}: MOND fails for galaxy clusters; framework doesn't address this
\end{enumerate}

\subsection{Comparison with Other Approaches}

\begin{table}[h]
\centering
\begin{tabular}{lccc}
\toprule
Approach & Derives $a_0$? & Explains cosmic coincidence? & Rigorous? \\
\midrule
Standard MOND & No (input) & No & N/A \\
Verlinde emergent gravity & Partially & Yes & Somewhat \\
$\mathbb{Z}_2$ framework & Formula & Yes & No \\
\bottomrule
\end{tabular}
\caption{Comparison of approaches to $a_0$}
\end{table}

\subsection{Honesty Assessment}

\textcolor{red}{\textbf{The $a_0 = cH_0/\mathbb{Z}_2$ formula is pattern-matching, not derivation.}}

What would constitute a real derivation:
\begin{enumerate}
    \item Start from the 7D orbifold action
    \item Show gravity modification emerges at low accelerations
    \item Derive the transition scale $a_0$ from geometry
    \item Explain why $\mathbb{Z}_2$ appears in the denominator
\end{enumerate}

\textbf{None of this has been done.}

\section{Conclusion}

The $T^3/\mathbb{Z}_2$ framework derives the MOND acceleration scale:
\begin{equation}
a_0 = \frac{cH_0}{\mathbb{Z}_2} = 1.18 \times 10^{-10} \text{ m/s}^2
\end{equation}

This explains the ``cosmic coincidence'' $a_0 \sim cH_0$ through shared geometric origin.

Key predictions:
\begin{enumerate}
    \item $a_0$ evolves with $H(z)$
    \item High-z galaxies show different MOND phenomenology
    \item Testable by JWST observations of $z > 5$ galaxies
\end{enumerate}

The derivation is plausible but incomplete---a rigorous holographic calculation would strengthen the result.

\section*{Acknowledgments}

The author thanks Stacy McGaugh for SPARC data and the MOND community for empirical calibrations.

\begin{thebibliography}{99}

\bibitem{milgrom1983}
Milgrom, M. (1983). ``A modification of the Newtonian dynamics.'' ApJ 270, 365.

\bibitem{mcgaugh2016}
McGaugh, S., Lelli, F., Schombert, J. (2016). ``Radial Acceleration Relation.'' Phys. Rev. Lett. 117, 201101.

\bibitem{sparc}
Lelli, F., McGaugh, S., Schombert, J. (2016). ``SPARC: Mass Models for 175 Disk Galaxies.'' AJ 152, 157.

\bibitem{verlinde2017}
Verlinde, E. (2017). ``Emergent Gravity and the Dark Universe.'' SciPost Phys. 2, 016.

\end{thebibliography}

\end{document}
